\documentclass[aspectratio=169,11pt]{beamer}
\usepackage[T1]{fontenc}
\usepackage[utf8]{inputenc}
\usepackage{graphicx}
\usepackage{xcolor}
\usepackage{hyperref}
\usepackage{xurl}
\usepackage{animate}

\hypersetup{colorlinks=true,linkcolor=blue,urlcolor=blue,citecolor=blue}

\definecolor{eifergreen}{RGB}{30,82,70}
\definecolor{eiferlight}{RGB}{238,245,241}
\definecolor{eiferdark}{RGB}{23,55,44}
\definecolor{eifergray}{RGB}{95,95,95}
\graphicspath{{../results/figures/}{results/figures/}{C:/Users/Hawi/Local codes/PV in Germany/results/figures/}}

\usetheme{default}
\setbeamertemplate{navigation symbols}{}
\setbeamertemplate{frametitle}{\vspace{0.20cm}\hspace{0.02\paperwidth}\insertframetitle\par\vspace{0.06cm}}
\setbeamercolor{frametitle}{fg=eiferdark}
\setbeamerfont{frametitle}{size=\Large,series=\bfseries}
\setbeamercolor{normal text}{fg=eiferdark,bg=white}
\setbeamerfont{normal text}{size=\small}
\setbeamersize{text margin left=0.012\paperwidth,text margin right=0.020\paperwidth}
\setbeamertemplate{footline}{}
\setbeamertemplate{headline}{}
\setbeamertemplate{frametitle continuation}{}

\newcommand{\compacttext}{\fontsize{6}{6.2}\selectfont\setlength{\parindent}{0pt}\setlength{\parskip}{0pt}\setlength{\parsep}{0pt}\setlength{\topsep}{0pt}\setlength{\partopsep}{0pt}\setlength{\itemsep}{0pt}\setlength{\leftmargini}{0.55em}\setlength{\labelsep}{0.15em}\setlength{\itemindent}{0pt}\setlength{\baselineskip}{6.2pt}\setlength{\lineskip}{0pt}\setlength{\lineskiplimit}{0pt}}
\newcommand{\source}[1]{\par\vspace{3mm}\begingroup\compacttext\begin{minipage}[t]{\linewidth}\noindent\textbf{Sources}\par #1\end{minipage}\endgroup}
\newcommand{\contentblock}[2]{%
\begingroup\compacttext\begin{minipage}[t]{\linewidth}
\if\relax\detokenize{#1}\relax\else
\noindent\textbf{Key messages}\par\let\itemize\enumerate\let\enditemize\endenumerate\setlength{\leftmargini}{1.3em}\setlength{\labelsep}{0.35em}#1\par\vspace{3mm}
\fi
\noindent\textbf{Methodology}\par #2\end{minipage}\endgroup}
\newcommand{\figslide}[5]{%
\begin{frame}{#1}
\begin{columns}[c,onlytextwidth]
\begin{column}{0.48\textwidth}\centering
\includegraphics[width=0.92\linewidth,height=0.65\textheight,keepaspectratio]{#2}
\end{column}
\begin{column}{0.50\textwidth}\raggedright\hspace{0pt}
\contentblock{#3}{#4}\source{#5}
\end{column}
\end{columns}
\end{frame}}

\begin{document}
\begin{frame}[plain]
\centering
\vfill
{\Huge\bfseries Photovoltaics in Germany\par}
\vspace{0.35cm}
{\Large Figures on capacity, density, electricity mix and redispatch\par}
\vfill
\end{frame}

\begin{frame}{Data sources}
\scriptsize
\begin{itemize}
\item \href{https://energy-charts.info/}{Energy-Charts.info}: electricity mix, installed capacity and district/state PV history.
\item \href{https://www.marktstammdatenregister.de/MaStR/}{Marktstammdatenregister}: operating-unit records, commissioning cohorts and EEG-exit inputs.
\item \href{https://www.smard.de/}{SMARD}: monthly redispatch observations and state aggregates.
\item \href{https://gdz.bkg.bund.de/index.php/default/digitale-geodaten/verwaltungsgebiete/verwaltungsgebiete-1-250-000-stand-31-12-vg250-31-12.html}{BKG VG250}: administrative boundaries and identifiers.
\item \href{https://www.destatis.de/DE/Themen/Branchen-Unternehmen/Landwirtschaft-Forstwirtschaft-Fischerei/Flaechennutzung/Tabellen/bodenflaeche-laender.html}{Destatis}: official state land-area totals.
\item \href{https://www.netzentwicklungsplan.de/}{Netzentwicklungsplan}: 2025 scenario tables and map pages.
\item \href{https://www.bundesnetzagentur.de/DE/Fachthemen/ElektrizitaetundGas/Versorgungssicherheit/Engpassmanagement/start.html}{BNetzA}: congestion-management and curtailment reports.
\end{itemize}
\end{frame}

\figslide{Electricity mix by source}{Figure_electricity-mix_germany_497-5-TWh.png}{\begin{itemize}\item Renewables supplied about 60\% of Germany's electricity generation in 2025.\item Wind and solar were the dominant renewable contributors in the mix.\end{itemize}}{Adapted from the Energy-Charts annual electricity-mix figure for Germany in 2025; the displayed categories and values are retained from the provider figure.}{\href{https://www.energy-charts.info/charts/energy_pie/chart.htm?l=en\&c=DE\&interval=year\&year=2025\&source=total\&legendItems=0wm}{Energy-Charts.info: annual electricity mix, Germany 2025}.}
\figslide{Installed electricity capacity by technology}{Figure_installed-power_germany_2026.png}{\begin{itemize}\item PV has become one of Germany's leading installed electricity-capacity technologies.\end{itemize}}{The Energy-Charts installed-power view is adapted for Germany 2026. Technology values are taken from the provider view in its capacity units; no forecast or generation conversion is applied.}{\href{https://energy-charts.info/charts/installed_power/chart.htm?l=en\&c=DE\&chartColumnSorting=default\&year=2026}{Energy-Charts.info: installed-power view}.}
\figslide{National PV density by district}{Figure_density_germany_labeled_2026.png}{\begin{itemize}\item High-capacity districts are not always the most intense once area is considered.\item The map shows where installed PV is concentrated spatially.\end{itemize}}{The district density values are taken from the Energy-Charts solar-per-area map and matched to BKG district geometries for the map and labels.}{\href{https://energy-charts.info/charts/installed_power_map/chart.htm?l=en\&c=DE\&source=solar}{Energy-Charts installed-power map}; \href{https://gdz.bkg.bund.de/index.php/default/digitale-geodaten/verwaltungsgebiete/verwaltungsgebiete-1-250-000-stand-31-12-vg250-31-12.html}{BKG VG250}.}

\begin{frame}{National PV density: 2000--2026}
\begin{columns}[c,onlytextwidth]
\begin{column}{0.48\textwidth}\centering
\begin{animateinline}[autoplay,controls,poster=last]{1}
\multiframe{27}{i=2000+1}{\includegraphics[width=0.92\linewidth,height=0.64\textheight,keepaspectratio]{C:/Users/Hawi/Local codes/PV in Germany/data/work/germany/Animation_density_germany_2000-2026/germany_pv_density_by_district_\i.png}}
\end{animateinline}
\end{column}
\begin{column}{0.50\textwidth}\raggedright\hspace{0pt}
\contentblock{\begin{itemize}\item PV expansion moves from a few concentrated areas toward broad national coverage.\item Southern and western districts remain prominent high-density areas as deployment spreads elsewhere.\item Growth in the total is accompanied by persistent spatial concentration.\end{itemize}}{One PNG frame is generated for each year from 2000--2026. District capacities are divided by district area, coloured with a fixed 0--780 kWp/km$^2$ scale, and embedded with \texttt{animateinline} and \texttt{multiframe}; year and total labels come from the workbook totals. The same yearly frames are also exported as a GIF for PowerPoint.}\source{\href{https://energy-charts.info/charts/installed_power_map/chart.htm?l=en\&c=DE\&source=solar}{Energy-Charts installed-power map}; \href{https://gdz.bkg.bund.de/index.php/default/digitale-geodaten/verwaltungsgebiete/verwaltungsgebiete-1-250-000-stand-31-12-vg250-31-12.html}{BKG VG250}.}
\end{column}
\end{columns}
\end{frame}

\figslide{East and West: observed comparison}{Figure_comparison_regions_east-vs-west_1999-2026.png}{\begin{itemize}\item East and West show a large difference in installed capacity, but a much smaller difference in capacity density.\item Normalising by area changes the regional comparison substantially.\end{itemize}}{Published district/state observations are assigned to fixed East and West groups, summed by year, and divided by summed group area. No interpolation or model adjustment is applied.}{\href{https://energy-charts.info/charts/installed_power_map/chart.htm?l=en\&c=DE\&source=solar}{Energy-Charts installed-power map}; \href{https://gdz.bkg.bund.de/index.php/default/digitale-geodaten/verwaltungsgebiete/verwaltungsgebiete-1-250-000-stand-31-12-vg250-31-12.html}{BKG VG250}; \href{https://www.destatis.de/DE/Themen/Branchen-Unternehmen/Landwirtschaft-Forstwirtschaft-Fischerei/Flaechennutzung/Tabellen/bodenflaeche-laender.html}{Destatis areas}.}
\figslide{North and South: observed comparison}{Figure_comparison_regions_north-vs-south_1999-2026.png}{\begin{itemize}\item The north--south difference in installed PV is less pronounced than the geographic impression suggests.\item Similar installed capacity does not imply similar generation: southern irradiance can raise output without changing the capacity comparison.\end{itemize}}{Published district/state observations are filtered into fixed North and South groups. Capacity is aggregated directly and density is calculated from total capacity divided by total group area; irradiance and generation are not modelled.}{\href{https://energy-charts.info/charts/installed_power_map/chart.htm?l=en\&c=DE\&source=solar}{Energy-Charts installed-power map}; \href{https://gdz.bkg.bund.de/index.php/default/digitale-geodaten/verwaltungsgebiete/verwaltungsgebiete-1-250-000-stand-31-12-vg250-31-12.html}{BKG VG250}; \href{https://www.destatis.de/DE/Themen/Branchen-Unternehmen/Landwirtschaft-Forstwirtschaft-Fischerei/Flaechennutzung/Tabellen/bodenflaeche-laender.html}{Destatis areas}.}
\figslide{East and West: capacity forecast}{Figure_capacity-and-density_regions_east-vs-west_1999-2045_NEP-ABC.png}{\begin{itemize}\item The East--West gap is much larger in total capacity than in capacity density.\item NEP means Network Development Plan: its A--C pathways are planning scenarios for future capacity, not a fitted prediction.\end{itemize}}{State observations are aggregated into fixed East and West groups, then divided by summed Destatis areas. Historical values come from Energy-Charts; the A--C band and B line are transcribed from NEP scenario tables.}{\href{https://energy-charts.info/charts/installed_power_map/chart.htm?l=en\&c=DE\&source=solar}{Energy-Charts installed-power map}; \href{https://www.netzentwicklungsplan.de/}{NEP 2037--2045, 2025 second draft}; \href{https://www.destatis.de/DE/Themen/Branchen-Unternehmen/Landwirtschaft-Forstwirtschaft-Fischerei/Flaechennutzung/Tabellen/bodenflaeche-laender.html}{Destatis state areas}.}
\figslide{North and South: capacity forecast}{Figure_capacity-and-density_regions_north-vs-south_1999-2045_NEP-ABC.png}{\begin{itemize}\item Observed North--South differences in installed PV are modest relative to the regional map impression.\item The NEP pathways imply different future expansion scales, not a measured current gap.\end{itemize}}{State observations are filtered into fixed North and South groups. Capacity is summed by year and density is recalculated from group capacity and area; future paths are copied from NEP scenario tables.}{\href{https://energy-charts.info/charts/installed_power_map/chart.htm?l=en\&c=DE\&source=solar}{Energy-Charts installed-power map}; \href{https://www.netzentwicklungsplan.de/}{NEP 2037--2045, 2025 second draft}; \href{https://www.destatis.de/DE/Themen/Branchen-Unternehmen/Landwirtschaft-Forstwirtschaft-Fischerei/Flaechennutzung/Tabellen/bodenflaeche-laender.html}{Destatis state areas}.}
\figslide{Baden-Württemberg versus Bavaria}{Figure_comparison_states_baden-wurttemberg-vs-bavaria_1999-2026.png}{\begin{itemize}\item Bavaria leads in both total capacity and capacity density; the density gap is smaller than the capacity gap.\item Both series show strong post-2010 expansion.\end{itemize}}{The two state records are selected from the Energy-Charts workbook. Capacity remains the published state total; density is calculated using each state's Destatis area, with no capacity normalisation.}{\href{https://energy-charts.info/charts/installed_power_map/chart.htm?l=en\&c=DE\&source=solar}{Energy-Charts installed-power map}; \href{https://gdz.bkg.bund.de/index.php/default/digitale-geodaten/verwaltungsgebiete/verwaltungsgebiete-1-250-000-stand-31-12-vg250-31-12.html}{BKG VG250}; \href{https://www.destatis.de/DE/Themen/Branchen-Unternehmen/Landwirtschaft-Forstwirtschaft-Fischerei/Flaechennutzung/Tabellen/bodenflaeche-laender.html}{Destatis state areas}.}
\figslide{Saxony-Anhalt versus Thuringia}{Figure_comparison_states_saxony-anhalt-vs-thuringia_1999-2026.png}{\begin{itemize}\item Saxony-Anhalt leads Thuringia in both total installed capacity and capacity density.\item The gap widens markedly after 2010, with Saxony-Anhalt reaching roughly twice Thuringia's 2026 capacity.\end{itemize}}{The two state identifiers are matched to the historical workbook and plotted with identical units, year coverage and density calculation. State capacity is retained; density is calculated from capacity and Destatis area.}{\href{https://energy-charts.info/charts/installed_power_map/chart.htm?l=en\&c=DE\&source=solar}{Energy-Charts installed-power map}; \href{https://gdz.bkg.bund.de/index.php/default/digitale-geodaten/verwaltungsgebiete/verwaltungsgebiete-1-250-000-stand-31-12-vg250-31-12.html}{BKG VG250}; \href{https://www.destatis.de/DE/Themen/Branchen-Unternehmen/Landwirtschaft-Forstwirtschaft-Fischerei/Flaechennutzung/Tabellen/bodenflaeche-laender.html}{Destatis state areas}.}
\figslide{Burgenland versus Saale-Holzland}{Figure_comparison_districts_burgenland-vs-saale-holzland_1999-2026.png}{\begin{itemize}\item The two districts diverge in both capacity and density.\item Local comparisons can reveal patterns hidden at state scale.\item District and state results should not be ranked as one population.\end{itemize}}{The district identifiers are matched to the Energy-Charts historical capacity and density records; BKG geometries provide the district map frame.}{\href{https://energy-charts.info/charts/installed_power_map/chart.htm?l=en\&c=DE\&source=solar}{Energy-Charts installed-power map}; \href{https://gdz.bkg.bund.de/index.php/default/digitale-geodaten/verwaltungsgebiete/verwaltungsgebiete-1-250-000-stand-31-12-vg250-31-12.html}{BKG VG250}.}
\figslide{Capacity versus density: all states}{Figure_capacity-vs-density_states_logx_2026.png}{\begin{itemize}\item The logarithmic capacity axis makes small and large states visible together.\item Capacity and density remain source-derived values; only the display axis is transformed.\item The scatter is descriptive; no regression is implied.\end{itemize}}{State capacity and area pairs are plotted with a logarithmic capacity coordinate. Only the coordinate transform changes; capacity, density and labels remain source-derived values.}{\href{https://energy-charts.info/charts/installed_power_map/chart.htm?l=en\&c=DE\&source=solar}{Energy-Charts installed-power map}; \href{https://www.destatis.de/DE/Themen/Branchen-Unternehmen/Landwirtschaft-Forstwirtschaft-Fischerei/Flaechennutzung/Tabellen/bodenflaeche-laender.html}{Destatis state areas}.}
\figslide{Capacity versus density: non-city states}{Figure_capacity-vs-density_states_excl-city-states_logx_2026.png}{\begin{itemize}\item Removing city states reveals the larger-area pattern.\item The logarithmic axis prevents large states from dominating visibility.\item Density itself is not transformed.\end{itemize}}{The state table is filtered to non-city states. Capacity is displayed on a logarithmic axis for legibility, while density remains the linear capacity-per-area calculation.}{\href{https://energy-charts.info/charts/installed_power_map/chart.htm?l=en\&c=DE\&source=solar}{Energy-Charts installed-power map}; \href{https://www.destatis.de/DE/Themen/Branchen-Unternehmen/Landwirtschaft-Forstwirtschaft-Fischerei/Flaechennutzung/Tabellen/bodenflaeche-laender.html}{Destatis state areas}.}
\figslide{Capacity versus density: city states}{Figure_capacity-vs-density_states_city-states_2026.png}{\begin{itemize}\item Small denominators amplify capacity density.\item High density does not mean high total capacity.\item City-state values need a different spatial interpretation.\end{itemize}}{The state capacity--density table is filtered to Berlin, Bremen and Hamburg after density is calculated from capacity and area. The subset is isolated to make the denominator effect visible.}{\href{https://energy-charts.info/charts/installed_power_map/chart.htm?l=en\&c=DE\&source=solar}{Energy-Charts installed-power map}; \href{https://www.destatis.de/DE/Themen/Branchen-Unternehmen/Landwirtschaft-Forstwirtschaft-Fischerei/Flaechennutzung/Tabellen/bodenflaeche-laender.html}{Destatis state areas}.}
\figslide{NEP B: ground-mounted and rooftop density}{Figure_ground-and-rooftop-density_germany_2024-2037-2045_NEP-B.png}{\begin{itemize}\item Rooftop and ground-mounted PV have different spatial signatures.\item The common scale makes the technologies comparable.\end{itemize}}{Six panels are reconstructed from NEP scenario-B map pages using one shared colour scale. District values are colour-read estimates rather than a published district table; BKG geometries provide the common map frame.}{\href{https://www.netzentwicklungsplan.de/}{NEP 2037--2045, 2025 first draft, scenario B}; \href{https://gdz.bkg.bund.de/index.php/default/digitale-geodaten/verwaltungsgebiete/verwaltungsgebiete-1-250-000-stand-31-12-vg250-31-12.html}{BKG VG250}.}
\figslide{Historical curtailment by state: grid congestion management (EinsMan)}{Figure_curtailment-einsman_germany_2016-2021.png}{\begin{itemize}\item Curtailment was strongly concentrated in the north: Schleswig-Holstein accounted for 72.3\% of the mapped total in 2016, while Lower Saxony accounted for 45.6\% in 2021.\item The concentration shifted within the north: Schleswig-Holstein's share fell to 32.0\% by 2021 as Lower Saxony became the largest contributor.\item The historical measure combines renewable plants with combined heat and power (CHP), so it is not a PV-only measure.\end{itemize}}{Here, EinsMan means grid congestion management: the temporary reduction of generation when network constraints require it. The figure transcribes BNetzA annual state values for renewable and CHP plants; map labels show each state's share of the national total. Published dashes remain missing, and no values are estimated from colour, neighbouring states or interpolation.}{\href{https://www.bundesnetzagentur.de/DE/Fachthemen/ElektrizitaetundGas/Versorgungssicherheit/Engpassmanagement/start.html}{BNetzA annual congestion-management reports, 2016--2021}.}
\figslide{Recent all-technology redispatch by state}{Figure_redispatch-reductions-gwh_germany_2023-2025.png}{\begin{itemize}\item The national technology mix is changing: wind accounted for about 51\% of the recorded total in 2023, versus 37\% in 2025; PV rose from about 4\% to 15\%.\item The geographic pattern is consistent with a hypothesis of more wind-related reductions in the north and more solar-related reductions in the south.\item This remains a hypothesis: the source does not cross technology and state.\end{itemize}}{The SMARD CSV is filtered to complete calendar years 2023--2025, grouped by state, and summed in GWh. This recent series covers all technologies. The national technology percentages are calculated separately from SMARD's technology table; the state table is not crossed with technology, so the north--south technology interpretation is not tested by this figure.}{\href{https://www.smard.de/resource/blob/217306/-/data-csv-data.csv}{SMARD monthly data CSV}; \href{https://www.smard.de/page/home/topic-article/211816/213384}{SMARD technology breakdown}; \href{https://www.smard.de/page/home/topic-article/211972/213254/netzengpassmanagement-je-bundesland-und-anrainerstaaten}{SMARD state redispatch information}.}
\figslide{Household price and PV selling tariff~\textcolor{red}{(ongoing)}}{Figure_selling-vs-consuming_germany_2010-2026.png}{\begin{itemize}\item Household electricity prices and PV selling tariffs move on different paths.\item The long-run tariff decline contrasts with the higher household-price series.\end{itemize}}{A BDEW household-price series is combined with selected EEG/PV tariff anchors on a common year axis. Only observed years through 2026 are shown; self-consumption, taxes, financing and operating costs are excluded.}{\href{https://www.bundesnetzagentur.de/DE/Fachthemen/ElektrizitaetundGas/ErneuerbareEnergien/EEG_Foerderung/Archiv_VergSaetze/start.html}{BNetzA/EEG tariff archive}; \href{https://www.bdew.de/service/daten-und-grafiken/bdew-strompreisanalyse/}{BDEW household-price series}.}
\figslide{PV capacity by MaStR unit-size band~\textcolor{red}{(ongoing)}}{Figure_additions-by-size_germany_2000-2025.png}{\begin{itemize}\item The sharp 2014 break is consistent with the EEG 2014 reform: lower support and a tighter deployment framework changed project economics.\item From 2015, larger ground-mounted projects also moved toward competitive support, reinforcing the change in the additions profile.\item The chart shows a policy/market break, not a change in PV technology.\end{itemize}}{The dated MaStR full export is filtered to operating PV units, grouped by commissioning year, and assigned to unit-size bands. Capacity is summed within each cohort; the register series shows the timing of the break but does not by itself establish its cause.}{\href{https://www.marktstammdatenregister.de/MaStR/}{Marktstammdatenregister full export}, revision 26.2, snapshot 13 October 2026; \href{https://www.bundesnetzagentur.de/DE/Fachthemen/ElektrizitaetundGas/ErneuerbareEnergien/EEG_Foerderung/Archiv_VergSaetze/artikel.html}{Bundesnetzagentur EEG tariff archive}; \href{https://www.bmwk.de/Redaktion/EN/FAQ/EEG-2017/fragen-und-antworten-zum-eeg-2017.html}{BMWE EEG support and auction framework}.}
\figslide{EEG-exit cohorts~\textcolor{red}{(ongoing)}}{Figure_eeg-exit_germany_2021-2037.png}{\begin{itemize}\item The figure tracks both the number of affected units and their installed capacity, so large systems matter even when they are few in number.\item The main exit wave arrives around 2030--2033, when early-2010s cohorts reach the end of support.\item Support expiry does not imply closure or repowering.\end{itemize}}{MaStR operating units are linked to commissioning records, grouped by commissioning year, and assigned an illustrative exit year equal to commissioning year plus 21. The two panels report cumulative units and capacity by cohort; the visual style is adapted from the PwC chart.}{\href{https://www.marktstammdatenregister.de/MaStR/}{Marktstammdatenregister operating-unit data}; EEG commissioning records; exit rule: commissioning year + 21; \href{https://www.pwc.de/de/energiewirtschaft/pwc-energyfacts-alte-photvoltaik-anlagen-foerderung.pdf}{PwC, \textit{Alte Photovoltaik-Anlagen: Ende der Förderung in Sicht}}.}

\end{document}
